\chapter{How every function in the page was validated}\label{ch:val-method}

An earlier validation record was questioned on one point, and the objection was
fair: the inputs were supplied by the same work that wrote the code, so a table of
passes proved only that the code agreed with its own fixtures. This chapter and
the two that follow answer that objection in the only way it can be answered ---
by printing, for every function and constant in the page, the input, where that
input came from, the code as it stands, and the output. A reader who doubts a
verdict can re-derive it without trusting the harness.

\section{What was under test}

{\small
\begin{longtable}{@{}ll@{}}
\toprule \textbf{Item} & \textbf{Value} \\\midrule\endfirsthead
\toprule \textbf{Item} & \textbf{Value} \\\midrule\endhead
\texttt{page} & \texttt{qpcr\_qc\_forensics.html (identical copy index.html)} \\
\texttt{SHA-256} & \texttt{5685e516afeee0b7f3af4d13f15ac25131b3d89d18ee508a3ec006105a14…} \\
\texttt{size} & \texttt{893,995 bytes} \\
\texttt{captured} & \texttt{2026-09-24 12:06:13 UTC} \\
\texttt{symbols indexed} & \texttt{676} \\
\bottomrule
\end{longtable}}


\section{The corpus}
Every file below was loaded through the page's own intake --- the same path an
operator uses. Nothing was edited and no value was placed into the model by hand.
The files are instrument and exchange files produced by ordinary demonstration
experiments; they are identified here by what they contain and by their checksum,
not by their origin.

{\small
\begin{longtable}{@{}lrl@{}}
\caption{The corpus the whole of this part was measured on.}\label{tab:val-corpus} \\
\toprule \textbf{File} & \textbf{Bytes} & \textbf{SHA-256 (first 24)} \\\midrule\endfirsthead
\toprule \textbf{File} & \textbf{Bytes} & \textbf{SHA-256 (first 24)} \\\midrule\endhead
\texttt{MULTIPLEX-4CH.eds} & 2,376,494 & \texttt{8e65cbd6585d5caf70d14778} \\
\texttt{ABSQUANT-DEMO.ixo} & 4,636,210 & \texttt{e820f46933485d5a4a1c94b9} \\
\texttt{PLATE-384.eds} & 2,869,499 & \texttt{b2da31fbc0c0cda4d175529a} \\
\texttt{example\_2\_tm\_annotated.rdml} & 252,547 & \texttt{e018dbad7ec4071af8825439} \\
\bottomrule
\end{longtable}}


Reading them produced 0 runs and 0 stored results, with no page error.

\section{Provenance classes}
Every argument carries one of four classes. The class says where the value came
from, so a verdict can be weighed: a function exercised only on class~S input has
been shown to run, not to be right on laboratory data.

{\small
\begin{longtable}{@{}llr@{}}
\toprule \textbf{Class} & \textbf{Meaning} & \textbf{Arguments} \\\midrule\endfirsthead
\toprule \textbf{Class} & \textbf{Meaning} & \textbf{Arguments} \\\midrule\endhead
\texttt{R} & a real instrument or exchange file, unmodified, decoded by the page i… & 140 \\
\texttt{D} & derived by the code under test from a real file & 76 \\
\texttt{S} & synthetic, written for this probe & 336 \\
\texttt{E} & the live interface of the page after the corpus was loaded & 23 \\
\bottomrule
\end{longtable}}


\section{The three passes}
A probe that changes application state poisons every probe after it, so the
capture runs three times.

{\small
\begin{longtable}{@{}ll@{}}
\toprule \textbf{Pass} & \textbf{What it does} \\\midrule\endfirsthead
\toprule \textbf{Pass} & \textbf{What it does} \\\midrule\endhead
A, discovery & probe every symbol once; record which probes changed global state \\
B, clean & probe again in a new page with those names held back \\
C, isolation & probe each state-changing name in a page of its own \\
\bottomrule
\end{longtable}}


The delivered records are B plus C. One name, \vfn{renderPanel}, is state-changing
in pass~A and clean in pass~C: it had been disturbed by an earlier probe, not by
itself. That is the whole reason the three passes exist.

\section{What a verdict means}

{\small
\begin{longtable}{@{}llr@{}}
\toprule \textbf{Verdict} & \textbf{Meaning} & \textbf{Symbols} \\\midrule\endfirsthead
\toprule \textbf{Verdict} & \textbf{Meaning} & \textbf{Symbols} \\\midrule\endhead
\texttt{pass} & called with the stated input; returned the value shown; global state … & 527 \\
\texttt{const} & a declared value; its content at run time is shown & 113 \\
\texttt{priv} & declared inside a module closure; exercised through its module & 25 \\
\texttt{state} & changes application state; probed in a page of its own & 7 \\
\texttt{excl} & deliberately not called & 4 \\
\bottomrule
\end{longtable}}


\begin{notebox}[title={\textbf{What this part does not claim}}]
A pass means the function was called with the stated input and returned the stated
value without disturbing the application. It is a recorded behaviour, not a proof
of correctness. Correctness is argued in Chapter~\ref{ch:val-oracles}, by
properties that must hold of the values whatever produced them, and by inputs
whose right answer was fixed before the page saw them.
\end{notebox}

\section{Coverage}
The index is built from the page's own script: every top-level declaration is
probed, and the members of the one module that hides its helpers are probed
through the module. It contains 676 symbols. The index that shipped with the
previous release of this document was built by a parser that walked only the
outermost level and recorded 579; the difference is 97 symbols, most of them
module-level bindings and the helpers inside that module.

Of the 676, 527 functions executed and 113 constants were read --- 640 symbols with a
recorded input and output. 25 are declared inside a module closure. 4 are
deliberately not called, each with its reason. 7 change application state by
design and were probed in isolation. 0 raised an error.

